\section{Conclusion}

This preregistered experiment examined how personalization affects delegation to algorithms in financial decision-making. We found that participants initially delegated more to an objective EV-maximizing algorithm than to a personalized algorithm that incorporated their individual preferences (66.9\% vs. 55.8\%, p = 0.0413). However, this difference did not persist in subsequent decisions (p = 0.4413). A large proportion of participants (75.1\%) never delegated to either algorithm.

Our findings challenge common assumptions about the benefits of personalization in AI systems. While personalization is often assumed to increase trust and adoption, our results show that it can actually deter delegation in certain contexts. Participants initially preferred an objective algorithm over a personalized one, even though the personalized algorithm should theoretically lead to higher expected utility for them.

This counterintuitive finding has important implications for the design of algorithmic advisors. Designers should carefully consider when to personalize recommendations and when to use objective algorithms. In contexts where objective criteria are available and valued, such as expected value maximization in financial decisions, objective algorithms may be more trusted and more likely to be used.

Our study also highlights the importance of user education and transparent communication about algorithm capabilities and limitations. Users may not fully understand the benefits of personalized algorithms, and designers should clearly communicate how these algorithms work and why they may be beneficial.

Future research should investigate the mechanisms underlying our finding that personalization deters delegation, test interventions to increase trust in personalized algorithms, and examine whether our findings generalize to other domains and cultural contexts. As AI continues to play an increasing role in decision-making, understanding when and why people delegate to algorithms will be crucial for designing systems that are both effective and trusted.

In conclusion, our study provides novel evidence that personalization can deter rather than encourage delegation to algorithms, challenging common assumptions about the benefits of personalization in AI systems. This finding contributes to the literature on algorithm appreciation, trust in AI, and human-AI collaboration, and has important practical implications for the design of algorithmic advisors in finance and other domains.